% results.tex
% Results (merged Experiments + Results):
%   setup preamble, then aggregate scorecard, preservation figures,
%   per-branch split, beta dependence.

We evaluated the three composers on a fixed set of per-asset JumpHMM
marginals calibrated against the real SPY growth-rate path. This
isolated the effect of the composition rule from the effect of
market-path randomness: all three composers consumed the same real SPY
history and the same per-ticker innovation draw, so metric differences
across composers reflected the composition rule alone. The ticker
universe was a $424$-asset US equity/ETF set drawn from Polygon.io
daily-close histories, filtered to tickers with a complete in-sample
trading history matching the reference length of AAPL
($T = 2{,}767$ daily closes covering January~2014 through
December~2024, the same window used by the companion
\texttt{JumpHMM.jl} validation paper); filtering on the common length
preserved $424$ tickers, of which the SPY total-return growth-rate
series was taken as the market path $g_m$ and the remaining $423$
tickers as non-market assets. For each non-market ticker $i$ we fitted
a \texttt{JumpHMM} marginal to the asset's own growth-rate history,
using the default configuration ($N = 100$ Laplace-partition states,
Student-$t$ emissions with $\nu = 5$, annualized growth rates via
\texttt{excess\_growth\_rates}, risk-free rate $r_f = 0$,
$\Delta t = 1/252$); jump parameters $(\epsilon, \lambda)$ were left at
the package defaults rather than tuned per ticker, which would improve
per-ticker marginal fit but was not required for a comparison that
tests the effect of the composition rule on top of whatever marginal
the generator produces. For each non-market ticker we then ran
ordinary least squares of the asset's growth-rate history on $g_m$,
yielding the calibrated quadruple
$(\alpha_i, \beta_i, R^2_{i,\real}, \sigma_{\eps,\real,i})$. The
universe spanned $\beta \in [0.01, 1.79]$ with median $\beta = 1.00$
and quartiles $(0.81, 1.00, 1.18)$; $R^2_{i,\real}$ ranged from
$0.001$ to $0.933$ (Fig.~\ref{fig:r2_distribution}), with only two
tickers above the $R^2_{\mathrm{preserve}} = 0.80$ threshold (QQQ at
$R^2 = 0.861$ and SPYG at $R^2 = 0.933$) and the highest-$R^2$
individual stock (BLK at $R^2 = 0.645$) a structural $\sim 0.22$
below; the branch assignment is therefore empirically
tracker-specific on this universe, not a threshold artifact.

The composition rules used a \emph{paired innovation} protocol
where applicable: for each non-market ticker and each replication
$r \in \{1, \dots, R\}$ we drew a single innovation path
$\teps_i^{(r)}$ of length $T$ from the ticker's JumpHMM marginal,
and the same $\teps_i^{(r)}$ was passed to the naive,
Gaussian-residual-SIM, and hybrid composers, so that differences
among those three reflect the composition rule and not generator
stochasticity. The \emph{Naive} composer set
$g_i^{(r)} = \alpha_i + \beta_i\, g_m + \teps_i^{(r)}$ with no
variance correction. The \emph{Gaussian SIM} baseline replaced the
JumpHMM draw with
$g_i^{(r)} = \alpha_i + \beta_i\, g_m + \sigma_{\eps,\real,i}\,\xi_i^{(r)}$,
$\xi_i^{(r)} \sim \mathcal{N}(0, I_T)$ drawn fresh per replication,
recovering the calibrated $R^2$ exactly but discarding any marginal
structure carried by $\teps$. The \emph{Hybrid} composer applied
Algorithm~\ref{alg:hybrid} with floor $f = 0.10$ and threshold
$R^2_{\mathrm{preserve}} = 0.80$. We evaluated three further
composers that draw their innovations from dedicated residual
generators, each fit to the real OLS residual series
$e_i(t) = G_i(t) - \alpha_i - \beta_i G_m(t)$. The
\emph{JumpHMM-on-residuals} composer fits a fresh JumpHMM marginal
to $e_i$ via a pseudo-price inversion (so the
\texttt{JumpHMM.jl}~\cite{alswaidanVarner2026jdiq} fit interface
that expects a price series accepts the residual target without
modification) and composes
$g_i^{(r)} = \alpha_i + \beta_i g_m + \teps_{\mathrm{resid},i}^{(r)}$
with no variance correction. The \emph{Block bootstrap} composer
draws a Politis--Romano stationary bootstrap
replicate~\cite{politisRomano1994} of the empirical residual series
with mean block length $\sqrt{T} \approx 50$ and composes
analogously. The \emph{GARCH(1,1)-$t$} composer fits a
conditional-variance model per ticker~\cite{bollerslev1986} and
substitutes a simulated residual path; $30$ of the $423$ tickers
yielded a non-stationary GARCH fit ($\alpha + \beta \ge 1$ at the
optimizer's termination) and were excluded from that composer only. All three composers skipped the
optional cross-sectional rank reorder in Algorithm~\ref{alg:hybrid} so
that per-ticker metrics reflected the marginal construction in
isolation; cross-sectional dependence is studied separately in the
companion \texttt{JumpHMM.jl} paper and is not the object of this
evaluation. For every composed series we computed three metric
families against the ticker's real growth-rate history and the real
SPY path: SIM recovery reported $(\hat\alpha_i, \hat\beta_i,
\hat R^2_i)$ from OLS of $g_i^{(r)}$ on $g_m$ along with the absolute
deviations $|\hat\beta - \beta|$ and $|\hat R^2 - R^2_{\real}|$;
marginal fidelity was assessed by Kolmogorov--Smirnov (KS) and
Anderson--Darling (AD) two-sample $p$-values against the real
marginal, Wasserstein-$1$ distance $W_1$, excess kurtosis $\kappa$,
and the classical Hill tail-index estimator
$\hat\xi = (k-1)^{-1}\sum_{i=1}^{k-1}\log(X_{(i)}/X_{(k)})$ on the
upper $5\%$ order statistics of $|g_i^{(r)}|$, so reported Hill values
are tail indices with $\hat\xi \approx 1/\alpha$ for a Pareto-tailed
distribution with shape $\alpha$; variance preservation was the ratio
$\mathrm{Var}(g_i^{(r)}) / \sigma_{\gen,i}^2$, using the real-data
variance as a proxy for the generator's targeted variance. Pass rates
were reported at $\alpha = 0.05$, and we reported $R = 100$
replications per ticker so that aggregate metrics were averages over
$423 \times 100 = 42{,}300$ per-ticker observations per composer;
the simulation seed was fixed at $1234$.

We began with the aggregate comparison across composers, computing per-asset medians of
each metric and then aggregating across the universe (Table~\ref{tab:aggregate}). 
Applying the protocol to the $423$ non-market tickers yielded one
composed series per $(\mathrm{ticker},\, \mathrm{composer},\,
\mathrm{replication})$ triple, for a total of
$423 \times 100 \times 3 = 126{,}900$ scored series.  All three composers recovered $\beta$ to
within a median absolute error of $0.018$ to $0.022$, consistent with
the identity $\mathrm{Cov}(g, g_m)/\sigma_m^2 \to \beta_i^{\eff}$
holding at all residual scales. The differences between composers
appeared everywhere else. The Gaussian SIM baseline recovered $R^2$
most tightly (median error $0.008$) because it targeted $R^2$ by
construction, but it failed marginal fidelity badly: KS pass rate
$0.6\%$, AD pass rate $0.5\%$, excess kurtosis $1.4$ (real data,
$\sim 13$), Hill tail index $0.22$ (versus $\sim 0.37$ on the
generator; Fig.~\ref{fig:tails}). The naive composition retained the generator marginal but
inflated its variance by $\beta_i^2 \sigma_m^2$, so its KS and AD
pass rates sat at $4.5\%$ and $2.7\%$. The hybrid composer held both
properties at once: KS pass rate $50.4\%$, AD pass rate $47.6\%$,
kurtosis $7.2$, Hill $0.36$; the $R^2$ recovery ($0.021$) came in
tighter than the naive baseline but looser than the Gaussian
baseline, reflecting that the variance-preserving branch targeted
$\sigma_{\gen,i}^2$ rather than $R^2_{i,\real}$.

The three residual-fit composers sidestep the variance-inflation
problem entirely by fitting a generator directly to $e_i$, which by
construction has variance $\sigma_{\eps,\real,i}^2$ and no
double-counted market component. Their aggregate KS pass rates were
$75.8\%$ (JumpHMM-on-residuals), $73.3\%$ (block bootstrap), and
$66.9\%$ (GARCH(1,1)-$t$), each higher than hybrid's $50.4\%$
(Table~\ref{tab:aggregate}, Fig.~\ref{fig:baseline}). The advantage
sat in the body of the distribution rather than the tails: hybrid's
median excess kurtosis ($7.17$) and Hill upper-tail index ($0.365$)
matched JumpHMM-on-residuals ($7.26$ and $0.365$) to within
sampling variability, confirming that the variance correction of
Eq.~\eqref{eq:scale} preserves the heavy-tailed character of the
full-return generator even when the rescaling is applied. Block
bootstrap and GARCH, by fitting their own residual generators,
reproduced the empirical residual series but with slightly lighter
tails ($\kappa = 6.68$ and $6.04$; Hill $0.330$ and $0.318$) than
either JumpHMM-based recipe.

The downstream consequence of the body-fidelity gap was the
question a per-ticker one-day Value-at-Risk backtest was designed to
answer (Table~\ref{tab:var}, Fig.~\ref{fig:var}). At $\alpha = 0.99$
the nominal exceedance rate is $1\%$; the naive composer
under-breached at $0.67\%$ and the Gaussian baseline over-breached
at $1.40\%$, both failing Kupiec unconditional coverage on the
majority of the universe ($45$ and $47\%$ pass rates). The hybrid,
JumpHMM-on-residuals, block bootstrap, and GARCH(1,1)-$t$ composers
all hit coverage within one-tenth of a percentage point of nominal
($0.95$, $0.99$, $1.09$, and $1.09\%$ respectively), and their
Kupiec pass rates clustered between $72$ and $86\%$. Hybrid and
JumpHMM-on-residuals were statistically indistinguishable on this
test: the $25$pp body-fidelity gap between them did not translate to
a Value-at-Risk accuracy difference. Hybrid's standalone advantage
is therefore scoped: it is the right rule when a full-return
per-asset generator is the given input, and it delivers tail
accuracy on par with a dedicated residual generator without
requiring a separate fitting stage.

To see where the KS advantage over the naive baseline came from, we
next plotted per-ticker KS pass rate and per-ticker marginal variance
ratio $\mathrm{Var}(g) / \sigma_{\gen}^2$ as functions of the
calibrated $\beta$ (Fig.~\ref{fig:preservation}). The variance-ratio
panel made the mechanism explicit: the naive scatter traced the
theoretical curve $1 + \rho_i$ with
$\rho_i = \beta_i^2 \sigma_m^2 / \sigma_{\gen,i}^2$, overshooting the
generator variance by $50\%$ to $75\%$ at high $\beta$, while hybrid
and Gaussian SIM sat on the unit line by design. The KS-pass-rate
panel showed the consequence: naive fell from a pass rate of
$\sim 0.25$ at $\beta \approx 0.3$ to below $0.01$ past
$\beta \approx 0.7$, while hybrid sat between $0.3$ and $0.9$ across
the full $\beta$ range. Gaussian SIM was pinned to zero throughout;
with its Gaussian residual it did not match the heavy-tailed real
marginal at any $\beta$. Having traced the KS advantage to variance
preservation, we next separated the two baselines' failure modes by
plotting per-ticker excess kurtosis and Hill tail index against
$\beta$ (Fig.~\ref{fig:tails}). The naive and hybrid composers
clustered together near the real-data kurtosis reference line
($\kappa_{\real} \approx 13$), with the difference between them
reflecting the variance inflation under naive rather than any
difference in tail shape. The Gaussian SIM series collapsed to
$\kappa \approx 1$ and Hill $\approx 0.22$, pinned below the generator
regardless of $\beta$; by replacing the heavy-tailed innovation with a
Gaussian draw the baseline removed the entire heavy-tail carrier. The
Hill panel confirmed the same separation at the tail: naive and
hybrid both reported $\sim 0.37$, while Gaussian SIM reported
$\sim 0.22$.

Having established aggregate preservation, we next examined the
branch-selection outcome on this universe and asked whether the
advantage held across the $\beta$ distribution. Applying the
branch-selection rule of Section~\ref{sec:method-algorithm} to the
$423$-ticker universe produced $421$ assignments to the
\texttt{hybrid} branch, $2$ to the \texttt{r2-preserve} branch, and
$0$ to the \texttt{hybrid-clipped} branch
(Fig.~\ref{fig:branch_map}, Table~\ref{tab:branch}). The two
$R^2$-preserving tickers were QQQ ($R^2 = 0.861$) and SPYG
($R^2 = 0.933$), both ETFs that track market-capitalization indices
closely; their KS pass rate under hybrid composition was $99\%$,
consistent with the $R^2$-preserving branch recovering both the
calibrated $\beta$ and the calibrated $R^2$ by construction. Clipping
did not trigger anywhere in the universe: the market loading ratio
$\rho_i = \beta_i^2 \sigma_m^2 / \sigma_{\gen,i}^2$ stayed below
$1 - f = 0.90$ for every ticker, because the JumpHMM marginals
carried enough variance that the market component never crowded out
the $f = 0.10$ idiosyncratic floor. To check whether the preservation
advantage held across the $\beta$ distribution, we bucketed the
universe by $\beta$ quartile and recomputed the aggregate metrics
within each bucket (Table~\ref{tab:beta_bucket}). The hybrid KS pass
rate declined gently from $67\%$ in the low-$\beta$ quartile to $44\%$
in the high-$\beta$ quartile; the Gaussian SIM pass rate stayed pinned
near zero across all buckets; the naive pass rate collapsed from
$16\%$ in Q1 to under $1\%$ in Q3 and Q4 as $\beta_i^2 \sigma_m^2$
grew. The gentle hybrid decline at high $\beta$ reflected the caveat
noted in Section~\ref{sec:method-algorithm}: as
$\beta^{\eff\,2}\sigma_m^2$ grew toward $(1 - f)\sigma_{\gen}^2$, the
market-driven component of $g$ tightened the marginal toward Gaussian
along its own axis, mildly attenuating the generator's tails in the
composed series. The effect was small enough that hybrid still beat
both baselines by more than an order of magnitude in every bucket.
